\section{Preprocessing}
\label{sec:preprocessing}

Real Antscan meshes (photogrammetry via \term{scAnt}~\cite{scant}, or
micro-CT) arrive with a wide range of defects: holes, disconnected
fragments, non-manifold edges and vertices, small debris, self-intersecting
and fused geometry, thin appendages that vanish under naïve simplification,
and highly variable pose and orientation. Two complementary repair
strategies were evaluated rather than assumed:

\begin{itemize}
\item \textbf{Alpha Wrap}~\cite{alphawrap} constructs an entirely new, watertight envelope
around the input geometry, regardless of how damaged the input is. Its
single scale parameter trades off two failure directions: too coarse an
envelope and thin anatomical structures (legs, antennae) disappear
into it; too fine, and the noise and defects the envelope was meant to
remove survive into the output.
\item \textbf{ManifoldPlus}~\cite{manifoldplus} instead repairs and regularises the input
surface directly, and was used as a fallback specifically for cases Alpha
Wrap's envelope construction did not adequately resolve.
\item \textbf{Conservative cleanup and reduction.} Before and after repair,
duplicate vertices, degenerate faces and loose vertices are removed and
normals recomputed, and topology-preserving quadric-error-metric (QEM)
reduction~\cite{qem} lowers face count without collapsing thin appendages.
\end{itemize}

\begin{investigation}{I12}{Does the rebuilt preprocessing pipeline actually
produce more reliably watertight, better-conditioned meshes?}
\invQuestion{Benchmarked against the existing pipeline on the same corpus,
does the Alpha-Wrap-plus-ManifoldPlus pipeline improve mesh topology?}
\invMethod{Watertightness and connected-component count across 44 processed
specimens, compared against the existing preprocessing pipeline; Alpha
Wrap's scale parameter swept toward progressively finer offsets.}
\invResult{$44$ of $44$ processed meshes were watertight, and $29$ of $44$
were reduced to a single connected component, both a substantial
improvement in topology over the existing pipeline. Some configurations at
the finest offsets encountered memory limitations at $32$\,GB.}
\invVerdict{ADOPT}{Retained as the production preprocessing pipeline.
Whether this translated into more accurate downstream registration is a
separate question, answered in \S\ref{sec:prep-null}: it did not, to a
measurable degree (\F{F2}): watertightness is a precondition for reliable
fitting, not, on its own, a source of anatomical accuracy.}
\end{investigation}

This preprocessing stage is wrapped in a production data-preparation module,
\allowbreak\texttt{prepare\_antscan\_\allowbreak data\_for\_mesh\_\allowbreak
fitting.py}, which converts raw
Antscan exports into a fitting-ready corpus: consistent scale and
orientation, organised clean/processed variants, and batch-ready input for
the registration stage proper. The registration stage itself, run through
\texttt{fitter\_3d/\allowbreak optimise.py}, preserves every intermediate optimisation
stage (\texttt{Stage\_0\_init}, \texttt{Stage\_1\_default},
\texttt{Stage\_2\_deform\_coarse}, \texttt{Stage\_3\_deform\_fine}) along
with per-specimen parameter files, reconstructed meshes, and loss curves,
so that any registration result in this report can be re-inspected or
re-run rather than taken on trust. This reproducibility infrastructure is
what made the controlled comparisons throughout Chapters~\ref{ch:structure}
and~\ref{ch:diagnostics} possible in the first place: a controlled variant
is only informative if the configuration that produced it is preserved
exactly.
